\section{Which route? Memory, holonomy and the present}
\label{sec:route}

A transfer can reach the same target state by different routes. We ask what
the target must remember about the route, and whether there is a present that
does not depend on it. The answers are simple, and we give them because a
criterion that ignores the route can fail in exactly these ways.

\paragraph{The current state need not be enough.}
Two histories can end in the same current state and still differ in what they
predict. Then no rule that reads only the current state can predict both
correctly. The observation is elementary.

\begin{proposition}[Current state is not always sufficient]
% claim: TTN-04
\label{prop:collision}
If two routes have the same current code and different futures, then no
function of the current code returns both futures. On the two-copy model there
are two such histories: the same current state, with different futures once
funding is taken into account.
\end{proposition}

What repairs the failure? If the pair of current code and route memory is
enough to predict the future, then either the current code alone is already
enough, or the route memory separates every pair of histories that share a
current code and differ in their futures. So when two histories collide in this way, a target needs information beyond
its current code, and route memory is one way to supply it. In a given domain,
a licence of the kind below shows that it is not needed.

\paragraph{When route memory may be dropped.}
Such an argument can take one of three standard forms
\citep{Newman1942, BaaderNipkow1998}. In each, the declared rewrites of routes
preserve the boundary and the executed actions; without that, a shared normal
form says nothing about endpoints. In summary:
\begin{itemize}
\item \emph{Normalization}: rewrites lower a rank, and every route reaches one
shared normal form.
\item \emph{Confluence}: any two routes can be rewritten to a common one.
\item \emph{Flatness}: each event acts by an invertible forward clock. Critical
pairs exist, cover all routes and are completed by valid completion rows.
Routes are connected under completion moves, and rewrites reach a shared normal
form. The completion rows satisfy their action equations, and the clock's
holonomy around each completion row is the identity.
\end{itemize}
The complete hypotheses are those of the Lean sufficiency theorems indexed in
Appendix~\ref{app:index}.

\begin{theorem}[Licences for dropping route memory]
% claim: TTH-09
\label{thm:licences}
Let a finite system of routes come with a comparison domain: routes with the
same boundary, admissible and closed under the declared rewrites. If the domain
carries a normalization, confluence or flatness licence, then any two routes in
the domain reach the same endpoint.
\end{theorem}

Each licence is sufficient, and we prove nothing more. In particular, a licence
is not necessary. On a fixed domain of three one-event routes, two routes have
equal endpoints and no licence holds. So equal endpoints alone do not show that
the route system is flat. The flatness licence we use is the stronger adopted
form: completion ranks and a shared unique normal form, not boundary equality
alone.

\paragraph{Holonomy is not an arrow.}
Path dependence around a loop, holonomy in the broad sense
\citep{Nakahara2003}, can look like a direction in time. It is not one.

\begin{proposition}[Holonomy without an arrow]
% claim: TTN-03
\label{prop:holonomy}
There is a finite route system with two reversible route maps that do not
commute, and whose record update has no strict step. It has nontrivial holonomy
and no endogenous arrow.
\end{proposition}

\paragraph{No present fixed by the boundary.}
Finally, suppose routes share a boundary, and we ask for a ``present'' that
depends only on that boundary. Holonomy can rule this out.

\begin{proposition}[No canonical present]
% claim: TTN-05
\label{prop:now}
There is a two-phase clock with a translation that moves the designated
present, phase zero, for which the local present along the identity route and
along the loop cannot be written as a function of the shared boundary.
\end{proposition}

\begin{figure}[tbp]
\centering
% F5: schematic of the fixed boundary and the mod-two translation action.
\begin{tikzpicture}[x=1cm,y=1cm,font=\small,
 state/.style={draw,rounded corners=2pt,fill=black!3,minimum width=3.5cm,minimum height=.85cm,align=center},
 route/.style={-{Stealth[length=2mm]},semithick}]
\node[state] (boundary) at (7,3.3) {Shared boundary};
\node[state] (early) at (2.6,.75) {Present at the initial phase};
\node[state] (late) at (11.4,.75) {Present at the translated phase};
\draw[route] (boundary.west) to[out=180,in=90] node[left,align=center,fill=white,inner sep=2pt] {identity\\route} (early.north);
\draw[route,densely dashed] (boundary.east) to[out=0,in=90] node[right,align=center,fill=white,inner sep=2pt] {loop\\route} (late.north);
\draw[route] (early.east) to[bend left=20] node[above] {clock translation} (late.west);
\draw[route] (late.west) to[bend left=20] node[below] {clock translation} (early.east);
\node[align=center] at (7,-.5) {Equal boundary, different local presents\\Nonzero holonomy; no boundary-only choice of present};
\node[font=\footnotesize] at (7,-1.25) {Schematic of the clock control};
\end{tikzpicture}

\caption{The clock control behind Proposition~\ref{prop:now}. Two routes share a boundary, but the loop translates the clock's phase, so the local presents differ and no present is fixed by the boundary alone.}
\label{fig:clock}
\end{figure}

The argument is one line: the two routes share the boundary, but their
presents differ. We do not claim more, and in particular nothing about routes
with genuinely different boundaries or about time as experienced. The lesson for
continuation is modest but definite. A target that inherits its present from a
transfer inherits it relative to a route. Keeping route memory is one way to
handle this. A normalization, confluence or flatness licence is a sufficient
reason to drop it, and it is not the only possible reason. A holonomy that
moves the designated present rules out a present fixed by the boundary alone.
A nontrivial transformation that leaves the designated present in place does
not.
